%%%%%%%%%%%%%%%%%%%%%%% CHAPTER - 5 %%%%%%%%%%%%%%%%%%%%\\
\chapter{Implementation and Module Description}
\label{C5} %%%%%%%%%%%%%%%%%%%%%%%%%%%%
\graphicspath{{images/}{Figures/}}
\noindent\rule{\linewidth}{2pt}
\textit{This Chapter presents the concrete software implementation and module-wise architectural decomposition of Campus Agent, detailing the full-stack technology stack, backend services, frontend single-page application, and specialized functional modules.}
%%%%%%%%%%%%%%%%%%%%%%%%%%%%%%%%%%%%%%%%%%%%%%%%%%%%%%%%%%%%%%%%%%%%%%%%%%%%%%%%%%
\section{Technology Stack Description} \label{S5.1}
The Campus Agent software platform is built using a modern, decoupled full-stack architecture optimized for high performance, responsiveness, and zero reliance on paid external cloud APIs:

\begin{itemize}
	\item \textbf{Backend Framework (Python \& FastAPI):} Selected for its asynchronous ASGI performance, automatic OpenAPI documentation, and native schema validation through Pydantic. Python provides extensive data manipulation and regular expression capabilities.
	\item \textbf{Frontend Application Framework (React 18 \& Vite):} React facilitates reactive, component-driven user interfaces. Vite serves as the modern build tool offering instant Hot Module Replacement (HMR) and optimized rollup production bundling.
	\item \textbf{Iconography \& Visual Elements (\texttt{lucide-react}):} Provides consistent, lightweight SVG icons across navigation, form fields, match badges, and Kanban cards.
	\item \textbf{Data Ingestion \& PDF Processing (\texttt{pypdf}):} Pure-Python library capable of in-memory binary PDF parsing, extracting text layers without external binary runtime dependencies.
	\item \textbf{Persistence Layer (JSON \& Browser \texttt{LocalStorage}):} File-based JSON datastore (\texttt{internships.json}) for opportunity cataloging; browser \texttt{LocalStorage} for persistent Kanban state tracking.
\end{itemize}

%%%%%%%%%%%%%%%%%%%%%%%%%%%%%%%%%%%%%%%%%%%%%%%%%%%%%%%%%%%%%%%%%%%%%%%%%%%%%%%%%%
\section{Backend Module Implementation} \label{S5.2}
The backend application resides in the \texttt{backend/} directory and is structured into five core Python modules:

\begin{enumerate}[label=\textbf{\arabic*.}]
	\item \textbf{\texttt{main.py} (FastAPI Server \& REST Routing):} Initialises the FastAPI application instance, configures permissive Cross-Origin Resource Sharing (CORS) middleware, and mounts REST endpoints:
	\begin{itemize}
		\item \texttt{POST /api/recommend}: Ingests candidate profile parameters and invokes the deterministic matching algorithm.
		\item \texttt{POST /api/parse-resume}: Receives multipart PDF uploads, extracting candidate credentials via \texttt{pypdf}.
		\item \texttt{POST /api/alert}: Triggers SMTP deadline alerts or records simulated dispatch logs.
		\item \texttt{POST /api/scrape}: Dispatches multi-agent scraping runs across specified enterprise portals.
		\item \texttt{POST /api/import-scraped}: Merges newly harvested opportunities into the primary database.
		\item \texttt{GET /api/health}: Returns service operational status and the total count of loaded opportunities.
	\end{itemize}
	\item \textbf{\texttt{agent.py} (Matching \& Reassessment Engine):} Implements the two-stage matching logic, duration unit normalization (weeks vs. months), calendar date to ordinal index conversion, score computation, and shortlist transition matrix generation.
	\item \textbf{\texttt{models.py} (Data Contract Schemas):} Formulates strongly-typed Pydantic classes ensuring validation of all incoming and outgoing payloads.
	\item \textbf{\texttt{scraper\_agent.py} (Web Scraping Agent):} Contains enterprise career URL generators, HTTP crawling agents, Groq LLM API integrations, and local NLP fallback logic.
	\item \textbf{\texttt{alerts.py} (Notification Engine):} Implements secure SMTP client connections and localized file logging in \texttt{data/last\_alert\_email.txt}.
\end{enumerate}

%%%%%%%%%%%%%%%%%%%%%%%%%%%%%%%%%%%%%%%%%%%%%%%%%%%%%%%%%%%%%%%%%%%%%%%%%%%%%%%%%%
\section{Frontend Single-Page Application} \label{S5.3}
The frontend is encapsulated in the \texttt{frontend/} directory, structured into modular React components coordinated by \texttt{App.jsx}:
\begin{itemize}
	\item \textbf{\texttt{App.jsx}:} Root application orchestrator managing global state, active tab navigation (Opportunity Finder, Kanban Tracker, Multi-Agent Scraper), What-If simulation parameters, and asynchronous recommendation dispatching.
	\item \textbf{\texttt{api.js}:} Asynchronous fetch wrapper connecting client UI components to FastAPI endpoints with structured error handling.
	\item \textbf{\texttt{index.css}:} Custom, modern stylesheet implementing responsive layouts, card elevations, color-coded match badges, and smooth transition animations.
\end{itemize}

%%%%%%%%%%%%%%%%%%%%%%%%%%%%%%%%%%%%%%%%%%%%%%%%%%%%%%%%%%%%%%%%%%%%%%%%%%%%%%%%%%
\section{Candidate Intake and PDF Resume Extraction Module} \label{S5.4}
The candidate intake interface (\texttt{StudentForm.jsx}) provides a dual-mode profile configuration interface:
\begin{itemize}
	\item \textbf{Manual Profile Configuration:} Provides text inputs and dropdown selectors for skills, domain interests, degree programmes, cohort years, CGPA, maximum duration constraints, and preferred start dates.
	\item \textbf{One-Click PDF Resume Dropzone:} Allows students to upload standard PDF resumes. Upon upload, \texttt{/api/parse-resume} streams the binary file, scans for 45+ technical keywords, identifies degree titles, extracts academic years, parses numeric CGPA metrics, and automatically populates the form fields in real time.
\end{itemize}

%%%%%%%%%%%%%%%%%%%%%%%%%%%%%%%%%%%%%%%%%%%%%%%%%%%%%%%%%%%%%%%%%%%%%%%%%%%%%%%%%%
\section{Ranked Recommendation and Decision Presentation Module} \label{S5.5}
The results presentation component (\texttt{ResultsList.jsx}) renders the outputs produced by \texttt{agent.py}:
\begin{itemize}
	\item \textbf{Featured Top Pick Card:} Displays the overall \#1 recommendation in a prominent banner, detailing the match level, matched skills, domain overlap, and an explainable rationale sentence verifying compliance with duration and availability dates.
	\item \textbf{Ranked Eligible Opportunities Grid:} Presents eligible positions sorted in descending order of relevance score, complete with badges indicating match tier (High, Medium, Low), required skills, duration, and application links.
	\item \textbf{Ineligible Opportunities Audit Trail:} Details all excluded opportunities with explicit, red-accented justification statements specifying violated constraints (e.g., "Duration of 12 weeks exceeds allowed 8 weeks" or "CGPA 8.2 is below required 8.5").
\end{itemize}

%%%%%%%%%%%%%%%%%%%%%%%%%%%%%%%%%%%%%%%%%%%%%%%%%%%%%%%%%%%%%%%%%%%%%%%%%%%%%%%%%%
\section{Interactive What-If Eligibility Parameter Simulator} \label{S5.6}
Integrated directly within \texttt{App.jsx}, the What-If Simulator enables candidates to dynamically explore how changes in their academic or scheduling profile impact their eligibility:
\begin{itemize}
	\item \textbf{Real-Time Parameter Sliders:} Candidates adjust CGPA, cohort year, maximum duration, and skill additions using responsive sliders and inputs.
	\item \textbf{Debounced Asynchronous Queries:} Slider changes trigger debounced API calls (400 ms delay), recalculating the entire recommendation pool dynamically without freezing the browser interface.
	\item \textbf{Impact Sensitivity Feedback:} Instantly reveals how elevating CGPA (e.g., from 8.2 to 8.65) re-qualifies previously excluded positions (e.g., unlocking \textit{Data Analytics Intern}).
\end{itemize}

%%%%%%%%%%%%%%%%%%%%%%%%%%%%%%%%%%%%%%%%%%%%%%%%%%%%%%%%%%%%%%%%%%%%%%%%%%%%%%%%%%
\section{Drag-and-Drop Application Lifecycle Kanban Tracker} \label{S5.7}
To support candidates through the entire job search lifecycle, Campus Agent provides a persistent Kanban application tracking pipeline:
\begin{itemize}
	\item \textbf{Four Pipeline Stages:} Opportunities move seamlessly across \textit{Interested}, \textit{Applied}, \textit{Interviewing}, and \textit{Offers} columns.
	\item \textbf{Browser Persistence:} Kanban states are serialized to JSON and persisted within the browser's \texttt{LocalStorage} (\texttt{campus\_agent\_tracked\_applications}), ensuring tracking records survive session restarts and page refreshes.
	\item \textbf{One-Click Bookmark Integration:} Any eligible recommendation card can be added directly to the Kanban board with a single click.
\end{itemize}

%%%%%%%%%%%%%%%%%%%%%%%%%%%%%%%%%%%%%%%%%%%%%%%%%%%%%%%%%%%%%%%%%%%%%%%%%%%%%%%%%%
\section{Multi-Agent Career Crawler and Importer Module} \label{S5.8}
The career harvester module (\texttt{CareerScraper.jsx} and \texttt{scraper\_agent.py}) extends the platform's reach beyond internal university notices:
\begin{itemize}
	\item \textbf{Enterprise Target Selector:} Allows users to select from 20+ top technology companies (Google, Microsoft, Meta, Amazon, Apple, Nvidia, Uber, etc.) or input custom job circular URLs.
	\item \textbf{Live Extraction Pipeline:} Executes multi-agent crawling, strips boilerplate DOM markup, and parses job parameters via Groq LLaMA-3.1 or the local NLP fallback engine.
	\item \textbf{Review \& Pool Import:} Displays harvested positions with title, organization, required skills, and deadlines, offering an administrative button to merge validated listings into \texttt{internships.json}.
\end{itemize}

%%%%%%%%%%%%%%%%%%%%%%%%%%%%%%%%%%%%%%%%%%%%%%%%%%%%%%%%%%%%%%%%%%%%%%%%%%%%%%%%%%
\section{Mock Technical Interview and Skill Gap Roadmap Module} \label{S5.9}
To actively assist students in technical interview preparation:
\begin{itemize}
	\item \textbf{Mock Technical Interview Engine:} Automatically generates technical verification questions based on the required skills of the candidate's top recommended role (e.g., core Python concepts, SQL queries, or machine learning architectures). The candidate enters their answers, and an automated evaluator grades their response out of 10.0 based on key conceptual keywords.
	\item \textbf{Curated Learning Roadmaps:} For skills required by an internship but absent from the student's profile, the system links students to curated, verified open-access educational roadmaps (e.g., Kaggle Learn, Fast.ai, MDN Web Docs, and AWS Skill Builder).
\end{itemize}

%%%%%%%%%%%%%%%%%%%%%%%%%%%%%%%%%%%%%%%%%%%%%%%%%%%%%%%%%%%%%%%%%%%%%%%%%%%%%%%%%%
\section{Summary} \label{S5.10}
This Chapter described the complete implementation architecture of Campus Agent, detailing the technological stack, backend REST endpoints, frontend React components, PDF resume parsing engine, What-If simulator, persistent Kanban tracker, multi-agent web scraper, and mock interview preparation tools.
